% Created 2025-02-17 lun 01:42
% Intended LaTeX compiler: pdflatex
\documentclass[11pt]{article}
\usepackage[utf8]{inputenc}
\usepackage[T1]{fontenc}
\usepackage{graphicx}
\usepackage{longtable}
\usepackage{wrapfig}
\usepackage{rotating}
\usepackage[normalem]{ulem}
\usepackage{amsmath}
\usepackage{amssymb}
\usepackage{capt-of}
\usepackage{hyperref}
\usepackage{minted}

\renewcommand{\contentsname}{Tabla de contenidos}
\usepackage{parskip}
\usepackage[utf8]{inputenc} %% For unicode chars
\usepackage{placeins}
\usepackage[margin=2.50cm]{geometry}
\usepackage[T1]{fontenc}
\usepackage{mathpazo}
\linespread{1.05}
\usepackage[scaled]{helvet}
\usepackage{courier}
\hypersetup{colorlinks=true,linkcolor=blue}
\RequirePackage{fancyvrb}
\DefineVerbatimEnvironment{verbatim}{Verbatim}{fontsize=\small,formatcom = {\color[rgb]{0.5,0,0}}}
\usepackage{mdframed}
\BeforeBeginEnvironment{minted}{\begin{mdframed}}
\AfterEndEnvironment{minted}{\end{mdframed}}
\setcounter{secnumdepth}{3}
\date{}
\title{Cliente Vue para API Rest con Laravel 10}
\hypersetup{
 pdfauthor={},
 pdftitle={Cliente Vue para API Rest con Laravel 10},
 pdfkeywords={},
 pdfsubject={},
 pdfcreator={Emacs 28.1 (Org mode 9.6.17)}, 
 pdflang={Spanish}}
\begin{document}

\maketitle
\setcounter{tocdepth}{3}
\tableofcontents

\setlength\parindent{10pt}
\section{Cliente en Vue.js para API REST de Laravel 10}
\label{sec:org304e2b2}

Este documento describe los pasos necesarios para crear un cliente en
Vue.js que consuma la API REST de Laravel 10 con autenticación JWT.

\subsection{El servidor}
\label{sec:orgd9bc374}
Primero hay que tener en cuenta que el servidor, después de hacer:ç
\begin{minted}[]{bash}
sail up -d
\end{minted}
tenemos que hacer
\begin{minted}[]{bash}
sail artisan serve
\end{minted}
\subsection{Instalación y Configuración del Proyecto Vue}
\label{sec:org2b886c4}

\subsubsection{1.1 Crear un nuevo proyecto Vue con Vite}
\label{sec:org3947a9d}
Ejecuta el siguiente comando para generar un proyecto Vue:

\begin{minted}[]{bash}
npm create vite@latest laravel-client --template vue
cd laravel-client
npm install
\end{minted}

\subsubsection{1.2 Instalar Axios y Vuex para manejo de estado}
\label{sec:org4daf60c}
Axios se usará para realizar peticiones HTTP y Vuex para manejar la autenticación:

\begin{minted}[]{bash}
npm install axios vuex@4
\end{minted}

\subsubsection{1.3 Configurar las variables de entorno}
\label{sec:org28618a6}
Crea un archivo \texttt{.env} en la raíz del proyecto Vue y agrega la URL de la API Laravel:

\begin{minted}[]{env}
VITE_API_URL=http://localhost/api/v1
\end{minted}
\subsection{Configurar App.vue}
\label{sec:org10b3947}
Configurar App.vue

Sustituye el contenido de src/App.vue por lo siguiente para integrar
Vue Router y mostrar las rutas correctamente:
\begin{minted}[]{html}
<template>
  <div>
    <nav>
      <router-link to="/">Home</router-link> |
      <router-link to="/login">Login</router-link>
      <router-link to="/create-post">Crear Post</router-link>
    </nav>
    <router-view />
  </div>
</template>

<script>
export default {
  name: "App",
};
</script>

<style>
nav {
  margin-bottom: 20px;
}
nav a {
  margin-right: 10px;
}
</style>

\end{minted}

\subsection{Configurar la Conexión con la API}
\label{sec:orge406a0b}

\subsubsection{2.1 Crear una instancia de Axios}
\label{sec:orga1963c4}
Crea un archivo \texttt{src/services/api.js}:

\begin{minted}[]{javascript}
import axios from 'axios';

const api = axios.create({
    baseURL: import.meta.env.VITE_API_URL,
    headers: {
        'Content-Type': 'application/json',
        'Accept': 'application/json',
    }
});

export default api;
\end{minted}

\subsection{Manejo de Autenticación con Vuex}
\label{sec:orged00ea2}

\subsubsection{3.1 Configurar el Store de Vuex}
\label{sec:org1351661}
Crea un archivo \texttt{src/store/index.js}:

\begin{minted}[]{javascript}
import { createStore } from 'vuex';
import api from '@/services/api';

export default createStore({
    state: {
        user: null,
        token: localStorage.getItem('token') || null
    },
    mutations: {
        setUser(state, user) {
            state.user = user;
        },
        setToken(state, token) {
            state.token = token;
            localStorage.setItem('token', token);
            api.defaults.headers['Authorization'] = `Bearer ${token}`;
        },
        logout(state) {
            state.user = null;
            state.token = null;
            localStorage.removeItem('token');
            delete api.defaults.headers['Authorization'];
        }
    },
    actions: {
        async login({ commit }, credentials) {
            const response = await api.post('/auth/login', credentials);
            commit('setToken', response.data.access_token);
            commit('setUser', await api.post('/auth/me'));
        },
        async logout({ commit }) {
            await api.post('/auth/logout');
            commit('logout');
        },
        async fetchUser({ commit }) {
            const response = await api.post('/auth/me');
            commit('setUser', response.data);
        }
    }
});
\end{minted}

\subsection{Crear Componentes para los Posts}
\label{sec:org8d42f0a}

\subsubsection{4.1 Componente de Lista de Posts}
\label{sec:orgddff21f}
Crea \texttt{src/components/PostList.vue}:

\begin{minted}[]{javascript}
<template>
  <div>
    <h1>Lista de Posts</h1>
    <transition name="fade">
      <div v-if="flashMessage" class="flash-message" @click="flashMessage = ''">
        {{ flashMessage }}
      </div>
    </transition>
    <ul>
      <li v-for="post in posts" :key="post.id">
        <h3>{{ post.title }}</h3>
        <p>{{ post.description }}</p>
        <img v-if="post.photo" :src="post.photo" alt="Imagen del post" class="post-image" />
      </li>
    </ul>
  </div>
</template>

<script>
import api from '@/services/api';

export default {
  data() {
    return {
      posts: [],
      flashMessage: localStorage.getItem('flashMessage') || ''
    };
  },
  async mounted() {
    try {
      const response = await api.get('/posts');
      this.posts = response.data.data;
      if (this.flashMessage) {
        setTimeout(() => {
          this.flashMessage = '';
          localStorage.removeItem('flashMessage');
        }, 3000);
      }
    } catch (error) {
      console.error('Error al cargar los posts', error);
    }
  }
};
</script>

<style>
.post-image {
  max-width: 100%;
  height: auto;
  margin-top: 10px;
}

.flash-message {
  background-color: #4caf50;
  color: white;
  padding: 10px;
  margin-bottom: 15px;
  text-align: center;
  border-radius: 5px;
  cursor: pointer;
}

.fade-enter-active, .fade-leave-active {
  transition: opacity 0.5s;
}
.fade-enter, .fade-leave-to {
  opacity: 0;
}
</style>
\end{minted}

\subsection{Agregar Vista para Crear un Post}
\label{sec:orgb93e7e2}

\subsubsection{4.1 Crear el Componente CreatePost.vue}
\label{sec:orgbfa8e00}
Crea un archivo src/components/CreatePost.vue:
\begin{minted}[]{java}
<template>
  <div>
    <h1>Crear Nuevo Post</h1>
    <form @submit.prevent="submitPost" enctype="multipart/form-data">
      <input v-model="title" type="text" placeholder="Título" required />
      <textarea v-model="description" placeholder="Descripción" required></textarea>
      <input type="file" @change="handleFileUpload" required />
      <button type="submit">Publicar</button>
      <p v-if="message" :class="{'error': isError, 'success': !isError}">{{ message }}</p>
    </form>
  </div>
</template>

<script>
import api from '@/services/api';

export default {
  data() {
    return {
      title: '',
      description: '',
      image: null,
      message: '',
      isError: false
    };
  },
  methods: {
    handleFileUpload(event) {
      this.image = event.target.files[0];
    },
 async submitPost() {
  try {
    const formData = new FormData();
    formData.append('title', this.title);
    formData.append('description', this.description);
    formData.append('image', this.image);

    const response = await api.post('/posts', formData, {
      headers: {
        'Content-Type': 'multipart/form-data'
      }
    });

    localStorage.setItem('flashMessage', response.data.message || 'Post creado con éxito!');
    this.$router.push('/');
  } catch (error) {
    console.error('Error al crear el post', error);
  }
}
  }
};
</script>

<style>
.error {
  color: red;
}
.success {
  color: green;
}
</style>
\end{minted}

\subsection{Configurar Rutas con Vue Router}
\label{sec:orgae13e10}

\subsubsection{5.1 Instalar Vue Router}
\label{sec:org1f978b1}
\begin{minted}[]{bash}
npm install vue-router@4
\end{minted}

\subsubsection{5.2 Configurar las Rutas}
\label{sec:org251d928}
Crea el archivo \texttt{src/router/index.js}:

\begin{minted}[]{javascript}
import { createRouter, createWebHistory } from 'vue-router';
import PostList from '@/components/PostList.vue';
import Login from '@/components/Login.vue';
import CreatePost from '@/components/CreatePost.vue';

const routes = [
    { path: '/', component: PostList },
    { path: '/login', component: Login }
    { path: '/create-post', component: CreatePost }
];

const router = createRouter({
    history: createWebHistory(),
    routes,
});

export default router;
\end{minted}

Modifica \texttt{src/main.js} para usar Vue Router y Vuex:

\begin{minted}[]{javascript}
import { createApp } from 'vue';
import App from './App.vue';
import router from './router';
import store from './store';

const app = createApp(App);
app.use(router);
app.use(store);
app.mount('#app');
\end{minted}

\subsection{Crear un Componente de Login}
\label{sec:org200732a}

Crea \texttt{src/components/Login.vue}:

\begin{minted}[]{javascript}
<template>
 <div>
  <h1>Iniciar Sesión</h1>
  <form @submit.prevent="handleLogin">
    <input v-model="email" type="email" placeholder="Correo" required>
    <input v-model="password" type="password" placeholder="Contraseña" required>
    <button type="submit">Ingresar</button>
  </form>
 </div>
</template>

<script>
import { mapActions } from 'vuex';

export default {
 data() {
     return {
         email: '',
         password: ''
     };
 },
 methods: {
   ...mapActions(['login']),
   async handleLogin() {
   try {
       await this.login({ email: this.email, password: this.password });
       this.$router.push('/');
   } catch (error) {
       console.error('Error en el login', error);
   }
   }
 }
};
</script>
\end{minted}

\subsection{Ejecutar el Proyecto}
\label{sec:org94fde5d}
Para iniciar el servidor Vue, ejecuta:

\begin{minted}[]{bash}
npm run dev
\end{minted}

Esto iniciará el frontend en \texttt{http://localhost:5174}.  Y el \texttt{80} para
 producción.

\subsection{Consideraciones para Nginx como Reverse Proxy}
\label{sec:orgdb0dc17}

Si el frontend Vue está en una máquina separada de Laravel detrás de
Nginx, asegúrate de que la configuración de Nginx permite solicitudes
CORS correctamente. Agrega esta configuración al bloque del servidor
en nginx.con
\begin{minted}[]{java}
location /api/ {
    proxy_pass http://backend-server:8000/api/;
    proxy_set_header Host $host;
    proxy_set_header X-Real-IP $remote_addr;
    proxy_set_header X-Forwarded-For $proxy_add_x_forwarded_for;
    proxy_set_header X-Forwarded-Proto $scheme;
}
\end{minted}
Si el frontend y backend están en el mismo servidor, asegúrate de que
el reverse proxy esté configurado para servir ambos correctamente.

Con esta actualización, el cliente Vue.js puede interactuar con el
backend Laravel detrás de Nginx sin problemas.

\subsection{Mejoras Futuras}
\label{sec:org1353b10}
\subsubsection{{\bfseries\sffamily DONE} Agregar un formulario para crear posts.}
\label{sec:orgcebbb47}
\subsubsection{Implementar edición y eliminación de posts.}
\label{sec:org6e9c7c7}
\subsubsection{{\bfseries\sffamily DONE} Agregar middleware en rutas para restringir acceso si no está autenticado.}
\label{sec:org511eac6}
----
\section{Usar Laravel Sail para Crear una Aplicación Front-end con Vue.js (Alternativa)}
\label{sec:org9073953}

Puedes usar Laravel Sail para crear tu aplicación front-end con
Vue.js. Aquí te dejo los pasos que puedes seguir para hacerlo:

\subsection{1. Inicia Laravel Sail}
\label{sec:org0790d75}

Asegúrate de que tu aplicación Laravel Sail esté en funcionamiento:

\begin{minted}[]{sh}
./vendor/bin/sail up -d
\end{minted}

\subsection{2. Crear la Aplicación Vue con Vite}
\label{sec:org1c50d51}

Utiliza el contenedor de Laravel Sail para crear la aplicación Vue.js
usando Vite:

\begin{minted}[]{sh}
./vendor/bin/sail npm create vite@latest laravel-client --template vue
\end{minted}

\subsection{3. Navegar al Directorio del Cliente Vue}
\label{sec:org08d44e5}

Después de crear la aplicación Vue, navega al directorio del proyecto:

\begin{minted}[]{sh}
cd laravel-client
\end{minted}

\subsection{4. Instalar Dependencias}
\label{sec:org4df0a3f}

Instala las dependencias necesarias utilizando npm dentro del
contenedor de Sail:

\begin{minted}[]{sh}
./vendor/bin/sail npm install
\end{minted}

\subsection{5. Configurar el Proxy en Vite}
\label{sec:org43667af}

Para que Vite pueda comunicarse con tu API Laravel, configura un proxy en \texttt{vite.config.js}:

\begin{minted}[]{js}
export default {
  server: {
    proxy: {
      '/api': 'http://localhost', // o la dirección donde se ejecuta tu aplicación Laravel
    },
  },
};
\end{minted}

\subsection{6. Ejecutar la Aplicación Vue}
\label{sec:org05c4f79}

Para ejecutar la aplicación Vue dentro del contenedor de Sail, usa el
siguiente comando:

\begin{minted}[]{sh}
./vendor/bin/sail npm run dev
\end{minted}

Esto debería iniciar el servidor de desarrollo de Vite y permitirte
ver tu aplicación Vue en tu navegador.

\subsection{Resumen de Comandos}
\label{sec:org786f40c}

\begin{minted}[]{sh}
# Iniciar Laravel Sail
./vendor/bin/sail up -d

# Crear la aplicación Vue
./vendor/bin/sail npm create vite@latest laravel-client --template vue

# Navegar al directorio del cliente Vue
cd laravel-client

# Instalar dependencias
./vendor/bin/sail npm install

# Configurar proxy en vite.config.js
# (Editar el archivo vite.config.js y agregar el proxy)

# Ejecutar la aplicación Vue
./vendor/bin/sail npm run dev
\end{minted}

Siguiendo estos pasos, podrás integrar una aplicación front-end Vue.js
con tu API REST Laravel utilizando Sail.

----
\section{Comparación entre Laravel Sail y Configuración Independiente}
\label{sec:orgeb017fe}

Ambas opciones tienen sus ventajas y desventajas. Aquí te dejo algunos
puntos a considerar para cada una:

\subsection{Usar Laravel Sail}
\label{sec:orgb3c008b}

\subsubsection{Ventajas}
\label{sec:org970f3fa}
\begin{itemize}
\item \textbf{\textbf{Entorno Consistente}}: Sail proporciona un entorno de desarrollo
consistente utilizando Docker, lo que puede reducir problemas de
configuración entre diferentes entornos de desarrollo.
\item \textbf{\textbf{Simplicidad}}: Configurar y administrar servicios como MySQL,
Redis, y otros es más sencillo con Sail, ya que todo está
preconfigurado.
\item \textbf{\textbf{Facilidad de Uso}}: Puedes aprovechar comandos simples de Sail
para ejecutar tareas comunes como la instalación de dependencias,
ejecutar pruebas, y más.
\end{itemize}

\subsubsection{Desventajas}
\label{sec:orge0f5ae6}
\begin{itemize}
\item \textbf{\textbf{Restricciones de Docker}}: Algunas veces, trabajar con Docker
puede introducir problemas específicos de contenedores, como temas
de rendimiento y configuración de red.
\item \textbf{\textbf{Curva de Aprendizaje}}: Si no estás familiarizado con Docker,
puede haber una pequeña curva de aprendizaje inicial.
\end{itemize}

\subsection{Hacerlo de Forma Independiente}
\label{sec:org70d710e}

\subsubsection{Ventajas}
\label{sec:org6d5e7bd}
\begin{itemize}
\item \textbf{\textbf{Flexibilidad}}: Puedes personalizar tu entorno de desarrollo según
tus necesidades específicas sin las restricciones de un contenedor
Docker.
\item \textbf{\textbf{Rendimiento}}: En algunos casos, un entorno no acoplado a Docker
puede tener un mejor rendimiento, especialmente en sistemas de
archivos y acceso a red.
\end{itemize}

\subsubsection{Desventajas}
\label{sec:orgf4aaf79}
\begin{itemize}
\item \textbf{\textbf{Configuración Manual}}: Necesitarás configurar manualmente
servicios como bases de datos, colas, etc., lo que puede ser tedioso
y propenso a errores.
\item \textbf{\textbf{Inconsistencias}}: Podrías encontrarte con problemas de
consistencia entre diferentes entornos de desarrollo y producción si
no tienes cuidado.
\end{itemize}

\subsection{Recomendación}
\label{sec:org5c7cac9}

Si ya estás utilizando Sail para tu API con Laravel, podría ser
conveniente seguir utilizando Sail para crear y gestionar tu
aplicación Vue, ya que te proporcionará un entorno de desarrollo
consistente y simplificará la configuración inicial. Usar Sail también
facilita la administración de dependencias y la integración continua.

Por otro lado, si prefieres un enfoque más flexible y personalizado, o
si ya tienes experiencia configurando entornos de desarrollo sin
Docker, podrías optar por crear la aplicación Vue de forma
independiente.

En resumen:
\begin{itemize}
\item \textbf{\textbf{Si la simplicidad y la consistencia son importantes para ti}},
sigue con Sail.
\item \textbf{\textbf{Si necesitas más flexibilidad y rendimiento}}, considera hacerlo
de forma independiente.
\end{itemize}
\end{document}
